\section{Two realizations of a modification}\label{sec:modifications}

We compare two descriptions of a line modification of a rank-$n$,
degree-one bundle. Framing the original bundle gives a projective
space with a division-algebra action; framing its slope-zero
modification gives a Drinfeld space with a linear-group action.
The common quotient will carry both degree-three classes. We then
apply the same modification to the universal rank-three-to-rank-two
surjection, obtaining the exact sequence that relates the two heights.

We work over $\Spd C$.  For a perfectoid test object $S$ over this
base, write $\mathcal X_S^{\mathrm{FF}}$ for its relative
Fargues--Fontaine curve, with untilt divisor
$i_\infty:\infty\hookrightarrow\mathcal X_S^{\mathrm{FF}}$.
At the fixed geometric point determined by $C$, write
$\mathcal X_C^{\mathrm{FF}}$ for the corresponding curve.
For a locally profinite group $A$, write
$B_CA=[\Spd C/\underline A]$ for its relative classifying stack;
all classifying maps in this section are over this base.
Write $\mathcal E_n^0=\cO(1/n)$, of rank $n$ and degree one.
A \emph{determinant lattice} of a slope-zero rank-$n$ bundle means a
$\Z_p$-lattice in its rank-one determinant local system.
Its frame group is
\[
 J_n=\{g\in\GL_n(\Q_p):|\det(g)|_p=1\}.
\]
This is the structure that will relate the two heights.

We use the following precise forms of the relative bundle theorems.
Let $\bar k$ be an algebraic closure of $k$.  Since $C^\flat$ is
algebraically closed, fix an extension $\bar k\hookrightarrow C^\flat$
of the embedding $k\hookrightarrow C^\flat$ underlying the chosen map
$\Spd C\to\Spd k$, and hence a compatible map
$\Spd C\to\Spd\bar k$.  Geometrically semistable bundles of slope zero
correspond to pro-\'etale $\Q_p$-local systems, by
\cite[Theorem II.2.19 and Corollary II.2.20]{FS}.
Pulling the geometrically trivial and basic strata of
\cite[Theorems III.2.4 and III.4.5]{FS} back along this map gives the
relative classifying descriptions $B_C\GL_n(\Q_p)$ and
$B_C\Aut(\mathcal E_n^0)$, respectively.
Thus their sheaves of frames are actual torsors in families.
The later arithmetic $G_F$-descent is established by the local
construction below.  We will apply the bundle statements after proving
the required slope conditions.  Geometric points alone are not being
used to identify arbitrary stacks.

\subsection{The lower and upper modifications}

\begin{lemma}\label{lem:lower-modification}
Let $\mathcal E_n$ be a form of $\mathcal E_n^0$, and let
$q:i_\infty^*\mathcal E_n\twoheadrightarrow\mathcal L$ be a quotient
line.  The elementary lower modification
\[
 0\longrightarrow\mathcal V_n\longrightarrow\mathcal E_n
 \longrightarrow i_{\infty *}\mathcal L\longrightarrow0
\]
is geometrically semistable of slope zero.  It therefore defines a
rank-$n$ rational local system.
\end{lemma}

\begin{proof}
The kernel is a vector bundle: locally at the Cartier divisor, the
quotient is a coordinate modulo a local equation of the divisor, and
the kernel has a basis obtained by multiplying that coordinate by the
local equation.  At a geometric point,
$\rank(\mathcal V_n)=n$ and $\deg(\mathcal V_n)=0$.
If it had a positive-slope Harder--Narasimhan piece, its maximal
destabilizing subbundle $\mathcal F$ would have rank $b<n$ and integral
degree $d\geq1$.  The injection $\mathcal F\hookrightarrow\mathcal E_n$
contradicts semistability of $\mathcal E_n$, because
\[
 \mu(\mathcal F)=\frac d b\geq\frac1b>\frac1n
       =\mu(\mathcal E_n).
\]
One may saturate its image in $\mathcal E_n$ without decreasing degree.
Hence $\mathcal V_n$ is semistable of degree zero at every geometric
point.  The relative slope-zero equivalence stated above now applies.
\end{proof}

Conversely, an elementary upper modification of $\cO^n$ at $\infty$
is determined by a line of modification directions.  Locally at a
geometric point, if $t$ is a uniformizer at $\infty$ and
$L_0=(B_{\mathrm{dR}}^+)^n$, its lattice is
\[
 L_\ell=L_0+B_{\mathrm{dR}}^+t^{-1}\widetilde\ell
       \subset t^{-1}L_0 ,
\]
where $\widetilde\ell$ lifts the direction line $\ell$.
Intrinsically the direction lies in
$(\cO(\infty)/\cO)^n$; tensoring every coordinate by the same
one-dimensional space identifies its projectivization with
$\mathbb P^{n-1}_C$.

\begin{lemma}\label{lem:drinfeld-locus}
The upper modification $\mathcal E_\ell$ is semistable of slope $1/n$
if and only if $\ell$ lies in no proper $\Q_p$-rational subspace.
Consequently its semistable locus is
\[
 \Omega_C^{n-1}
 =\mathbb P_C^{n-1}\setminus
       \bigcup_{\substack{W\subset\Q_p^n\\ \dim W=n-1}}\mathbb P(W_C).
\]
On this locus $\mathcal E_\ell$ is a form of $\mathcal E_n^0$.
\end{lemma}

\begin{proof}
If $\ell\subset W_C$ for a proper rational subspace $W$, modify
$W\otimes_{\Q_p}\cO$ along $\ell$.  The resulting rank-$\dim W$,
degree-one bundle embeds as a subbundle of $\mathcal E_\ell$; the
quotient is the unchanged trivial bundle associated to $\Q_p^n/W$.
Its slope is greater than $1/n$, so $\mathcal E_\ell$ is not semistable.

For the converse, take a saturated destabilizing subbundle
$\mathcal F\subset\mathcal E_\ell$ of rank $b<n$.
Its integral degree $d$ is at least one.  Set
$\mathcal F'=\mathcal F\cap\cO^n$.  The quotient
$\mathcal F/\mathcal F'$ embeds in $\mathcal E_\ell/\cO^n$ and has
length $\epsilon\leq1$.  Since $\mathcal F'$ embeds in $\cO^n$,
semistability of the latter gives
\[
 0\geq\deg(\mathcal F')=d-\epsilon\geq0.
\]
Thus $d=\epsilon=1$ and $\deg(\mathcal F')=0$.
The saturation of $\mathcal F'$ in $\cO^n$ cannot increase its degree:
a positive increase would contradict the same semistability inequality.
It is therefore already saturated.  All its subbundles have
nonpositive degree, so it is semistable of slope zero.  The
classification of bundles at a geometric point identifies it with
$W\otimes_{\Q_p}\cO$ for a $b$-dimensional rational subspace $W$;
here an inclusion between trivial bundles is a matrix over
$H^0(\mathcal X_C^{\mathrm{FF}},\cO)=\Q_p$.
The equality $\epsilon=1$ in the displayed local lattice description
then says precisely that $\ell\subset W_C$.

This proves the criterion at every geometric point.  The classification
in \cite[Theorem II.2.14]{FS} gives $\cO(1/n)$ for a semistable
rank-$n$, degree-one bundle, and the
relative basic-stratum theorem gives the claimed family of forms.
\end{proof}

\subsection{A common torsor with a determinant lattice}

The lower and upper modification criteria identify the same geometric
objects. A full lattice in the slope-zero local system would impose
extra data on these objects. We retain only its determinant lattice,
which is already supplied by the $P_n$-structure, and construct the
common space of frames with that condition.

Put
\[
 \Delta_n=\det(\mathcal E_n^0)(-\infty),\qquad
 V(\Delta_n)=H^0(\mathcal X_C^{\mathrm{FF}},\Delta_n).
\]
The latter is a one-dimensional $\Q_p$-space.  Choose a lattice
$\Lambda_n\subset V(\Delta_n)$.
The determinant action of
$\Aut(\mathcal E_n^0)=D_n^\times$ is the reduced norm.
The stabilizer of $\Lambda_n$ is
\[
 \{g\in D_n^\times:v_p(N_n(g))=0\}
       =\cO_{D_n}^{\times}=P_n.
\]
Accordingly a $P_n$-structure on a form $\mathcal E_n$ transports
$\Lambda_n$ to a determinant lattice on
$\det(\mathcal E_n)(-\infty)$.

These choices can be made with compatible arithmetic Galois action.
Take $F=W(k)[1/p]$, or a finite unramified extension of it.
The fixed bundle is obtained functorially from the Dieudonn\'e module
of $\Gamma_n$ over $k$, by extension to the period rings and Frobenius
descent; see \cite[Lemma 2.3.3 and the following paragraph,
pp.~23--24]{BMR2026}.  Thus it is defined on perfectoid test objects
over $\Spd F$.  Every endomorphism of $\Gamma_n$ is defined over $k$,
so this descent commutes with $P_n$.  The untilt divisor also descends:
its ideal is the kernel of the functorial untilt map $\theta$.
It follows that $\Delta_n$ descends over $\Spd F$.

Its geometric slope is zero.  Applying
\cite[Corollary II.2.20]{FS} on perfectoid test objects and descending
the resulting local systems therefore makes $V(\Delta_n)$ a
continuous rank-one representation of
$G_F=\Gal(\overline{\Q}_p/F)$.
Continuity is part of this conclusion: the descent automorphisms are
sections of the pro-\'etale $\Q_p^\times$-sheaf, hence continuous
functions on the profinite Galois parameters.  The coefficient sheaf
here is the one identified by \cite[Proposition II.2.5(ii)]{FS}.
The image of its valuation character in $\Z$ is a compact subgroup,
hence zero.  Every lattice $\Lambda_n$ is therefore $G_F$-stable.
We need no invariant choice of basis.  The cyclotomic character of
$G_F$ still has open image.

Let $\mathcal Z_n$ be the v-sheaf of injections
\[
 u:\cO^n\hookrightarrow\mathcal E_n^0
\]
whose cokernel is a line supported at $\infty$ and for which
\[
 \det(u):\det(\cO^n)\xrightarrow{\ \sim\ }\Delta_n
\]
takes the standard determinant lattice onto $\Lambda_n$.
The commuting actions of $P_n$ and $J_n$ are postcomposition and
inverse precomposition.
Galois acts on a map $u$ by transport along the semilinear descent of
its source and target.  On $\cO^n$ this descent fixes the standard
rational basis.  It follows directly that $\gamma(guj^{-1})
=g\gamma(u)j^{-1}$ for $\gamma\in G_F$, $g\in P_n$, and $j\in J_n$.
The determinant condition is preserved by stability of $\Lambda_n$.

\begin{proposition}\label{prop:two-realizations}
Let $\mathcal M_n$ classify a form of $\mathcal E_n^0$ with
$P_n$-structure and a quotient line at $\infty$.  There are
$G_F$-equivariant equivalences of v-stacks
\begin{equation}\label{eq:two-realizations}
 \mathcal M_n
 \simeq [\mathbb P(i_\infty^*\mathcal E_n^0)/P_n]
 \simeq [\mathcal Z_n/(P_n\times J_n)]
 \simeq [\Omega_C^{n-1}/J_n].
\end{equation}
Its framing torsor defines $\mathcal M_n\to B_CP_n$, and its universal
lower modification defines the indicated map $\mathcal M_n\to B_CJ_n$.
\end{proposition}

\begin{proof}
Fixing a quotient line of $i_\infty^*\mathcal E_n^0$ fixes its kernel
$\mathcal V_n$.  By \cref{lem:lower-modification} its frames exist
pro-\'etale locally.  Requiring a frame to preserve the determinant
lattice cuts its $\GL_n(\Q_p)$-torsor down to a $J_n$-torsor.
Thus
$\mathcal Z_n/J_n=\mathbb P(i_\infty^*\mathcal E_n^0)$
as v-sheaves.

In the other direction, a framed lower bundle gives an upper
modification $\cO^n\hookrightarrow\mathcal E_\ell$.
By \cref{lem:drinfeld-locus} its allowed directions are exactly
$\Omega_C^{n-1}$.  Its determinant lattice is induced by the standard
one on $\cO^n$.  The isomorphisms
$\mathcal E_\ell\simeq\mathcal E_n^0$ form a $D_n^\times$-torsor,
by the relative basic-stratum theorem.  The norm valuation
$D_n^\times\to\Z$ is surjective, so this torsor has, v-locally,
isomorphisms carrying the induced lattice onto $\Lambda_n$.
Those isomorphisms form precisely a $P_n$-torsor.
It follows that $\mathcal Z_n/P_n=\Omega_C^{n-1}$.

Both identifications commute with base change and the remaining group
action.  Taking the further quotients proves
\cref{eq:two-realizations}.  They are Galois compatible because every
operation used---forming the kernel, determinant, lattice condition,
and sheaf of frames---is compatible with the descent just constructed.
More explicitly, an upper direction is a line in
$\Q_p^n\otimes_{\Q_p}(\cO(\infty)/\cO)|_\infty$.
The second factor is common to all coordinates, so its semilinear
Galois scalar cancels upon projectivization.  The action on
$\Omega_C^{n-1}$ is consequently the standard arithmetic action,
commuting with $J_n$.  The lower frames form a Galois-equivariant
$J_n$-torsor, which gives the equivariant map to $B_CJ_n$.
Here Galois acts on the $\Spd C$ base of $B_CJ_n$ and trivially on
the group $J_n$; the analogous statement holds for $B_CP_n$.
This assertion does not require
the torsor itself to have a Galois-fixed frame.
\end{proof}

The two frame-forgetting maps can be pictured as follows; the upper
arrows are torsors for the indicated groups.
\[
\begin{tikzcd}[column sep=large,row sep=large]
 & \mathcal Z_n \arrow[dl,"J_n"'] \arrow[dr,"P_n"] & \\
 \mathbb P(i_\infty^*\mathcal E_n^0)
       \arrow[dr] && \Omega_C^{n-1}\arrow[dl]\\
 & \mathcal M_n &
\end{tikzcd}
\]
Only a determinant lattice was retained.  This explains the appearance
of $J_n$: its defining condition preserves the determinant lattice,
while the group as a whole preserves no full lattice in $\Q_p^n$.

\subsection{The filtration supplied by the two-height tower}

We have compared the two realizations at each rank. The remaining
construction relates those ranks using the universal surjection.

Return to the universal surjection
$f:\mathcal E_3\twoheadrightarrow\mathcal E_2$ over $Y_C$ from
\cref{prop:tower}; its all-surjections interpretation is
\cite[Proposition 2.7.1(3)]{BMR2026}.
Let
\[
 \pi:\mathcal T=\mathbb P(i_\infty^*\mathcal E_2)\longrightarrow Y_C
\]
parametrize quotient lines.  The linear $P_2$-descent on
$i_\infty^*\mathcal E_2$ gives an actual vector bundle on $Y_C$.
In particular $\mathcal T$ has its global tautological
$\cO_{\mathcal T}(1)$.

\begin{proposition}\label{prop:filtered-modification}
The tautological quotient of $\mathcal E_2$ and its composite with $f$
give maps $m_n:\mathcal T\to\mathcal M_n$ and an exact sequence
\begin{equation}\label{eq:filtered-modification}
 0\longrightarrow\ker(f)\longrightarrow\mathcal V_3
       \longrightarrow\mathcal V_2\longrightarrow0 .
\end{equation}
These are slope-zero bundles.  Their corresponding rational local
systems form an exact sequence of ranks $1,3,2$.
With the induced determinant lattices this sequence is classified by
a map $d:\mathcal T\to B_CD$, where
\[
 D=\left\{
 \begin{pmatrix}a&v\\0&A\end{pmatrix}:
       a\in\Z_p^\times,\quad A\in J_2,\quad v\in\Q_p^{1\times2}
       \right\}
   \simeq\Q_p^2\rtimes(\Z_p^\times\times J_2).
\]
The classifying maps of $\mathcal V_3$ and $\mathcal V_2$ are the
composites of $d$ with the relative maps induced by the inclusion
$\rho_3:D\hookrightarrow J_3$ and projection $\rho_2:D\to J_2$.
The $P_n$-classifying map of $m_n$ is the composite
$\mathcal T\xrightarrow{\pi}Y_C\to B_CP_n$ induced by
$q_n:G\to P_n$.
\end{proposition}

\begin{proof}
If $q:\mathcal E_2\twoheadrightarrow i_{\infty *}\mathcal L$
is the universal quotient, then
$\mathcal V_2=\ker(q)$ and $\mathcal V_3=\ker(qf)$.
Thus $\mathcal V_3=f^{-1}\mathcal V_2$, proving the exact sequence.
The two lower modifications are slope zero by
\cref{lem:lower-modification}.  The kernel of $f$ has rank one and
degree $\deg(\mathcal E_3)-\deg(\mathcal E_2)=0$, so it too has slope
zero.

To check exactness on local systems, pass to a pro-\'etale cover
trivializing the three slope-zero bundles.  Full faithfulness of the
relative correspondence identifies the maps with matrices of
continuous $\Q_p$-valued functions.  The matrices have the indicated
ranks on every geometric fiber; on the open loci of invertible minors
they give exact, locally split sequences of rational local systems.
These loci cover the base, proving exactness before descent.

Transport $\Lambda_3$ and $\Lambda_2$ to the determinant local systems
of $\mathcal V_3$ and $\mathcal V_2$.  The determinant isomorphism in
\cref{eq:filtered-modification} defines the lattice on the kernel as
their tensor quotient.  Locally choose a lattice-compatible basis of
the kernel and determinant-compatible frames of the quotient, and
lift the latter to the middle local system.  Such adapted frames form
a torsor; their change-of-frame matrices are exactly the displayed
matrices in $D$.  The two classifying maps follow by forgetting the
filtration or taking its quotient.  The construction of $m_n$ retains
the original $P_n$-structured form $\mathcal E_n$ on $Y_C$ and only
adds its quotient line.  Its $P_n$-framing torsor is therefore the
pullback of the associated framing torsor on $Y_C$, proving the last
assertion as a statement of actual torsors.
\end{proof}
